\documentclass[aspectratio=169]{beamer}

% ── Theme & colors ──────────────────────────────────────────────
\usetheme{Madrid}
\usecolortheme{default}
\definecolor{azuloscuro}{RGB}{10,40,80}
\definecolor{azulmedio}{RGB}{25,90,160}
\definecolor{acento}{RGB}{255,165,0}
\definecolor{fondoclaro}{RGB}{240,248,255}

\setbeamercolor{palette primary}{bg=azuloscuro,fg=white}
\setbeamercolor{palette secondary}{bg=azulmedio,fg=white}
\setbeamercolor{palette tertiary}{bg=azuloscuro,fg=white}
\setbeamercolor{palette quaternary}{bg=azuloscuro,fg=white}
\setbeamercolor{structure}{fg=azulmedio}
\setbeamercolor{frametitle}{bg=azuloscuro,fg=white}
\setbeamercolor{title}{fg=white}
\setbeamercolor{block title}{bg=azulmedio,fg=white}
\setbeamercolor{block body}{bg=fondoclaro,fg=black}
\setbeamercolor{alerted text}{fg=acento}
\setbeamertemplate{navigation symbols}{}
\setbeamertemplate{itemize items}[circle]

% ── Packages ────────────────────────────────────────────────────
\usepackage[utf8]{inputenc}
\usepackage[T1]{fontenc}
\usepackage[spanish]{babel}
\usepackage{amsmath,amssymb,bm}
\usepackage{tikz}
\usetikzlibrary{positioning}
\usetikzlibrary{arrows.meta,calc,decorations.markings,angles,quotes,3d,shapes.geometric}
\usepackage{xcolor}
\usepackage{booktabs}
\usepackage{multicol}

% ── Macros ──────────────────────────────────────────────────────
\newcommand{\vL}{\bm{L}}
\newcommand{\vr}{\bm{r}}
\newcommand{\vp}{\bm{p}}
\newcommand{\vtau}{\bm{\tau}}
\newcommand{\vomega}{\bm{\omega}}
\newcommand{\valpha}{\bm{\alpha}}
\newcommand{\vF}{\bm{F}}

% ── Title info ───────────────────────────────────────────────────
\title[Dinámica del Cuerpo Rígido]{\textbf{Dinámica de un Cuerpo Rígido}}
\subtitle{Rotación, Inercia y Movimiento Angular}
\author{Ricardo Leon Martinez}
\institute{DEMAT}
\date{19 de mayo de 2026}

% ════════════════════════════════════════════════════════════════
\begin{document}
% ════════════════════════════════════════════════════════════════

% ── PORTADA ─────────────────────────────────────────────────────
{
\setbeamertemplate{background}{
  \begin{tikzpicture}[remember picture,overlay]
    \fill[azuloscuro] (current page.south west) rectangle (current page.north east);
    % decorative arcs suggesting rotation
    \foreach \r/\op in {3/0.08,4.5/0.06,6/0.04,7.5/0.03}{
      \draw[white,opacity=\op,line width=1.2pt]
        (current page.center) ++(-2,0) circle(\r);
    }
    \draw[acento,line width=2pt,opacity=0.5]
      (current page.center) ++(-2,0) circle(2.2);
    % axis arrow
    \draw[-{Stealth[length=8pt]},acento,line width=2pt,opacity=0.7]
      ($(current page.center)+(-2,-3)$) -- ($(current page.center)+(-2,3.2)$);
    \node[white,opacity=0.15,font=\fontsize{120}{120}\selectfont] at
      ($(current page.center)+(4,-0.5)$) {$\omega$};
  \end{tikzpicture}
}
\begin{frame}[plain]
  \vspace{1.5cm}
  \begin{center}
    {\color{acento}\rule{0.55\textwidth}{1.5pt}}\\[0.4cm]
    {\color{white}\fontsize{22}{26}\selectfont\textbf{Dinámica de un Cuerpo Rígido}}\\[0.2cm]
    {\color{white!70}\large Rotación, Inercia y Movimiento Angular}\\[0.2cm]
    {\color{acento}\rule{0.55\textwidth}{1.5pt}}\\[0.8cm]
    {\color{white}\large Ricardo Leon Martinez}\\[0.1cm]
    {\color{white!60}\normalsize 19 de mayo de 2026}
  \end{center}
\end{frame}
}

% ── ÍNDICE ──────────────────────────────────────────────────────
\begin{frame}{Contenido}
  \begin{columns}[c]
    \column{0.5\textwidth}
    \tableofcontents
    \column{0.45\textwidth}
    \begin{tikzpicture}[scale=0.85]
      % Spinning disk diagram
      \fill[azulmedio!20] (0,0) ellipse(1.8 and 0.45);
      \fill[azulmedio!40] (0,0.15) ellipse(1.8 and 0.45);
      \draw[azuloscuro,line width=1.2pt] (0,0.15) ellipse(1.8 and 0.45);
      \draw[azuloscuro,line width=1.2pt] (-1.8,0.15) -- (-1.8,-0.15);
      \draw[azuloscuro,line width=1.2pt] (1.8,0.15) -- (1.8,-0.15);
      \draw[azuloscuro!50,line width=0.6pt] (0,0.15) ellipse(0.9 and 0.22);
      % axis
      \draw[acento,line width=1.5pt,dashed] (0,-1) -- (0,2.2);
      \draw[-{Stealth},acento,line width=1.5pt] (0,1.5) -- (0,2.2);
      % omega arc
      \draw[-{Stealth},azulmedio,line width=1.2pt]
        (0.4,0.55) arc[start angle=40,end angle=140,radius=0.6]
        node[midway,above,font=\small,azulmedio]{$\omega$};
      % labels
      \node[font=\small,azuloscuro] at (0,-0.6) {Disco rígido};
    \end{tikzpicture}
  \end{columns}
\end{frame}

% ── 1. CUERPO RÍGIDO ────────────────────────────────────────────
\section{Cuerpo Rígido}
\begin{frame}{¿Qué es un Cuerpo Rígido?}
  \begin{columns}[c]
    \column{0.48\textwidth}
    \begin{block}{Definición}
      Sistema de partículas donde la \alert{distancia entre cualquier par} de
      puntos permanece \textbf{constante} durante el movimiento.
    \end{block}
    \vspace{0.4cm}
    \begin{itemize}
      \item Traslación del CM
      \item \alert{Rotación} alrededor del CM
    \end{itemize}
    \vspace{0.3cm}
    \[
      \boxed{E_{\text{total}} = E_{\text{trans}} + E_{\text{rot}}}
    \]

    \column{0.48\textwidth}
    \begin{tikzpicture}[scale=0.9]
      % rigid body outline
      \fill[azulmedio!15,rounded corners=8pt]
        (-1.5,-1) -- (1.5,-0.8) -- (1.8,1) -- (-1.2,1.2) -- cycle;
      \draw[azuloscuro,line width=1.2pt,rounded corners=8pt]
        (-1.5,-1) -- (1.5,-0.8) -- (1.8,1) -- (-1.2,1.2) -- cycle;
      % particles
      \foreach \x/\y in {-0.8/-0.3, 0.2/0.4, 0.9/-0.2, -0.3/0.7, 0.5/0.8}{
        \fill[azulmedio] (\x,\y) circle(2.5pt);
      }
      % rigid bonds
      \draw[azulmedio!50,dashed,line width=0.7pt]
        (-0.8,-0.3)--(0.2,0.4)--(0.9,-0.2)--(-0.8,-0.3);
      \draw[azulmedio!50,dashed,line width=0.7pt]
        (0.2,0.4)--(-0.3,0.7)--(0.5,0.8)--(0.2,0.4);
      % CM
      \fill[acento] (0.1,0.2) circle(3.5pt);
      \node[font=\footnotesize,acento,right=3pt] at (0.1,0.2) {CM};
      % arrows
      \draw[-{Stealth},acento,line width=1.2pt] (0.1,0.2)--++(0.9,-0.3)
        node[right,font=\footnotesize]{$\bm{v}_{\text{cm}}$};
      \draw[-{Stealth},azulmedio,line width=1.2pt]
        (0.1,0.2)++(0.5,0.5) arc[start angle=20,end angle=100,radius=0.5]
        node[above,font=\footnotesize]{$\omega$};
    \end{tikzpicture}
  \end{columns}
\end{frame}

% ── 2. MOMENTO DE INERCIA ────────────────────────────────────────
\section{Momento de Inercia}
\begin{frame}{Momento de Inercia}
  \begin{columns}[c]
    \column{0.45\textwidth}
    \[
      \boxed{I = \sum_i m_i r_i^2 = \int r^2\, dm}
    \]
    \vspace{0.2cm}
    \begin{block}{Teorema de Steiner}
      \[
        I = I_{\text{cm}} + M d^2
      \]
    \end{block}
    \vspace{0.2cm}
    \small
    \begin{tabular}{lc}
      \toprule
      \textbf{Cuerpo} & $I$ \\
      \midrule
      Disco sólido   & $\dfrac{1}{2}MR^2$ \\[6pt]
      Esfera sólida  & $\dfrac{2}{5}MR^2$ \\[6pt]
      Varilla (CM)   & $\dfrac{1}{12}ML^2$ \\
      \bottomrule
    \end{tabular}

    \column{0.50\textwidth}
    \begin{tikzpicture}[scale=1.0]
      % Disk with particles at radius r
      \fill[azulmedio!15] (0,0) circle(1.8);
      \draw[azuloscuro,line width=1pt] (0,0) circle(1.8);
      % axis
      \fill[azuloscuro!60] (-0.08,-1.9) rectangle (0.08,1.9);
      \fill[acento] (0,0) circle(3.5pt) node[above right,font=\small]{$O$};
      % example particle
      \fill[azulmedio] (1.2,1.1) circle(4pt) node[right,font=\small]{$m_i$};
      \draw[acento,line width=1pt,dashed] (0,0)--(1.2,1.1)
        node[midway,above,sloped,font=\small]{$r_i$};
      % rotation arrow
      \draw[-{Stealth},azulmedio!80,line width=1.2pt]
        (1.9,0.3) arc[start angle=9,end angle=70,radius=1.9]
        node[above right,font=\small]{$\omega$};
      \node[font=\small,azuloscuro] at (0,-2.2){Disco $I=\tfrac{1}{2}MR^2$};
    \end{tikzpicture}
  \end{columns}
\end{frame}

% ── 3. TORQUE ───────────────────────────────────────────────────
\section{Torque}
\begin{frame}{Torque (Momento de Fuerza)}
  \begin{columns}[c]
    \column{0.48\textwidth}
    \[
      \boxed{\vtau = \vr \times \vF}
    \]
    \[
      |\vtau| = r\,F\sin\theta
    \]
    \vspace{0.2cm}
    \begin{itemize}
      \item Análogo rotacional de la \alert{fuerza}
      \item Dirección: regla de la mano derecha
      \item Unidades: $\text{N·m}$
    \end{itemize}

    \column{0.48\textwidth}
    \begin{tikzpicture}[scale=1.0]
      % pivot
      \fill[azuloscuro] (0,0) circle(4pt) node[below left,font=\small]{$O$};
      % arm r
      \draw[azulmedio,line width=1.8pt,-{Stealth}]
        (0,0) -- (2.5,0.8) node[above,font=\small]{$\vr$};
      % force F
      \draw[acento,line width=1.8pt,-{Stealth}]
        (2.5,0.8) -- (3.2,2.4) node[right,font=\small]{$\vF$};
      % angle arc
      \draw[gray,line width=0.8pt]
        (2.5,0.8)++(0.55,0.18) arc[start angle=17.7,end angle=75,radius=0.58]
        node[midway,right,font=\footnotesize]{$\theta$};
      % tau arc (torque)
      \draw[-{Stealth},azulmedio!60,line width=1.3pt]
        (0.7,-0.2) arc[start angle=-16,end angle=90,radius=0.73]
        node[above left,font=\small,azulmedio]{$\vtau$};
      % perpendicular line
      \draw[gray,dashed,line width=0.7pt]
        (2.5,0.8) -- (2.5,-0.5) node[below,font=\footnotesize]{$r\sin\theta$};
      \draw[gray,line width=0.7pt]
        (2.5,0) -- (2.65,0) -- (2.65,-0.15);
    \end{tikzpicture}
  \end{columns}
\end{frame}

% ── 4. MOMENTO ANGULAR ──────────────────────────────────────────
\section{Momento Angular}
\begin{frame}{Momento Angular}
  \begin{columns}[c]
    \column{0.48\textwidth}
    \textbf{Partícula:}
    \[
      \vL = \vr \times \vp = \vr \times m\bm{v}
    \]
    \textbf{Cuerpo rígido:}
    \[
      \boxed{\vL = I\,\vomega}
    \]
    \vspace{0.2cm}
    \begin{alertblock}{Conservación}
      Si $\vtau_{\text{neto}} = 0$:
      \[
        \vL = \text{cte}
      \]
    \end{alertblock}

    \column{0.48\textwidth}
    \begin{tikzpicture}[scale=0.88]
      % rotating body
      \fill[azulmedio!20,rounded corners=6pt]
        (-1.2,-0.5) rectangle (1.2,0.5);
      \draw[azuloscuro,line width=1pt,rounded corners=6pt]
        (-1.2,-0.5) rectangle (1.2,0.5);
      % axis through center
      \draw[gray,dashed,line width=0.8pt] (0,-1.6) -- (0,2.5);
      % omega
      \draw[-{Stealth},azulmedio,line width=1.5pt]
        (0.5,0.7) arc[start angle=30,end angle=150,radius=0.58]
        node[midway,above,font=\small]{$\omega$};
      % L vector
      \draw[-{Stealth},acento,line width=2pt]
        (0,0.5) -- (0,2.2) node[right,font=\small]{$\vL = I\vomega$};
      % r and p for a particle
      \fill[azulmedio] (1.0,0) circle(3pt);
      \draw[-{Stealth[length=5pt]},azulmedio!60,line width=0.9pt]
        (0,0)--(1.0,0) node[midway,below,font=\tiny]{$r$};
      \draw[-{Stealth[length=5pt]},acento!70,line width=0.9pt]
        (1.0,0)--(1.0,0.8) node[right,font=\tiny]{$\bm{p}$};
    \end{tikzpicture}
  \end{columns}
\end{frame}

% ── 5. ECUACIÓN DEL MOVIMIENTO ──────────────────────────────────
\section{Ecuación del Movimiento Rotacional}
\begin{frame}{Ecuación del Movimiento Rotacional}
  \begin{center}
    \begin{tikzpicture}
      \node[draw=azulmedio,rounded corners=8pt,fill=fondoclaro,
            inner sep=12pt,line width=1.5pt] {
        \Large$\displaystyle\vtau_{\text{neto}} = I\,\valpha$
      };
    \end{tikzpicture}
  \end{center}
  \vspace{0.1cm}
  \begin{columns}[c]
    \column{0.45\textwidth}
    \textbf{Analogías Traslación $\leftrightarrow$ Rotación}\\[0.3cm]
    \begin{tabular}{cc}
      \toprule
      \textbf{Traslación} & \textbf{Rotación} \\
      \midrule
      $m$           & $I$ \\
      $\bm{v}$      & $\vomega$ \\
      $\bm{a}$      & $\valpha$ \\
      $\bm{F}$      & $\vtau$ \\
      $\bm{F}=m\bm{a}$ & $\vtau=I\valpha$ \\
      \bottomrule
    \end{tabular}

    \column{0.50\textwidth}
    \begin{tikzpicture}[scale=0.9]
      % pulley
      \fill[azulmedio!20] (0,0) circle(1.2);
      \draw[azuloscuro,line width=1.2pt] (0,0) circle(1.2);
      \fill[acento] (0,0) circle(4pt);
      % rope
      \draw[azuloscuro,line width=1.5pt] (1.2,0) -- (1.2,-1.4);
      % hanging mass
      \fill[azulmedio!50] (0.75,-1.4) rectangle (1.65,-2.0);
      \draw[azuloscuro,line width=1pt] (0.75,-1.4) rectangle (1.65,-2.0);
      \node[font=\small] at (1.2,-1.7) {$M$};
      % torque arc
      \draw[-{Stealth},acento,line width=1.2pt]
        (0.3,1.1) arc[start angle=75,end angle=5,radius=1.15]
        node[right,font=\small]{$\tau$};
      % alpha label
      \node[azulmedio,font=\small] at (-0.6,0.3){$\alpha$};
      \draw[-{Stealth},azulmedio,line width=1pt]
        (-0.3,0.8) arc[start angle=100,end angle=170,radius=0.85];
      % weight
      \draw[-{Stealth},acento!80,line width=1pt]
        (1.2,-2.0) -- (1.2,-2.7) node[right,font=\small]{$Mg$};
    \end{tikzpicture}
  \end{columns}
\end{frame}

% ── 6. ENERGÍA ROTACIONAL ───────────────────────────────────────
\section{Energía Rotacional}
\begin{frame}{Energía Rotacional}
  \begin{columns}[c]
    \column{0.48\textwidth}
    \[
      \boxed{E_{\text{rot}} = \frac{1}{2}\,I\,\omega^2}
    \]
    \vspace{0.2cm}
    \textbf{Rodadura sin deslizamiento:}
    \[
      E_{\text{total}} = \underbrace{\frac{1}{2}mv_{\text{cm}}^2}_{\text{traslación}}
                       + \underbrace{\frac{1}{2}I_{\text{cm}}\omega^2}_{\text{rotación}}
    \]
    \vspace{0.2cm}
    \begin{block}{Trabajo rotacional}
      \[W = \int \tau\, d\theta\]
    \end{block}

    \column{0.48\textwidth}
    \begin{tikzpicture}[scale=1.0]
      % inclined plane
      \fill[gray!20] (0,0) -- (3.5,0) -- (3.5,-0.15) -- (0,-0.15) -- cycle;
      \draw[gray,line width=0.8pt] (0,0) -- (3.5,0);
      \fill[gray!30] (0,0) -- (3.5,0) -- (0,2.0) -- cycle;
      \draw[azuloscuro,line width=1pt] (0,0) -- (3.5,0) -- (0,2.0) -- cycle;
      % angle
      \draw[gray] (0.6,0) arc[start angle=0,end angle=29.7,radius=0.6]
        node[midway,right,font=\tiny]{$\theta$};
      % rolling disk
      \fill[azulmedio!30] (2.1,0.82) circle(0.42);
      \draw[azulmedio,line width=1.2pt] (2.1,0.82) circle(0.42);
      \draw[azuloscuro,line width=0.8pt] (2.1,0.82)--++(0.3,0.29);
      \draw[azuloscuro,line width=0.8pt] (2.1,0.82)--++(-0.3,-0.29);
      % v arrow
      \draw[-{Stealth},acento,line width=1.2pt]
        (2.1,0.82) -- (2.6,1.07) node[above,font=\small]{$v$};
      % omega arc
      \draw[-{Stealth},azulmedio!80,line width=1pt]
        (2.53,0.62) arc[start angle=-30,end angle=60,radius=0.42]
        node[right,font=\tiny]{$\omega$};
      % height h
      \draw[<->,acento!70,line width=0.7pt]
        (-0.2,0) -- (-0.2,2.0) node[midway,left,font=\small]{$h$};
    \end{tikzpicture}
  \end{columns}
\end{frame}

% ── 7. RESUMEN ──────────────────────────────────────────────────
\begin{frame}{Resumen}
  \begin{center}
  \begin{tikzpicture}[
    node distance=0.6cm and 1.0cm,
    box/.style={draw=azulmedio,fill=fondoclaro,rounded corners=6pt,
                minimum width=3.0cm,minimum height=0.9cm,
                font=\small,align=center,line width=1pt},
    formula/.style={font=\small\color{azuloscuro},align=center}
  ]
    \node[box] (I) {Momento de\\Inercia};
    \node[formula,right=0.4cm of I] (fI) {$I=\int r^2 dm$};

    \node[box,below=0.5cm of I] (tau) {Torque};
    \node[formula,right=0.4cm of tau] (ftau) {$\tau=rF\sin\theta$};

    \node[box,below=0.5cm of tau] (eqmov) {Ec.\ Movimiento};
    \node[formula,right=0.4cm of eqmov] (feq) {$\vtau=I\valpha$};

    \node[box,below=0.5cm of eqmov] (L) {Momento Angular};
    \node[formula,right=0.4cm of L] (fL) {$\vL=I\vomega$};

    \node[box,below=0.5cm of L] (E) {Energía Rot.};
    \node[formula,right=0.4cm of E] (fE) {$E_{\text{rot}}=\tfrac{1}{2}I\omega^2$};

    % arrows
    \draw[-{Stealth},acento,line width=1pt] (I.south)--(tau.north);
    \draw[-{Stealth},acento,line width=1pt] (tau.south)--(eqmov.north);
    \draw[-{Stealth},acento,line width=1pt] (eqmov.south)--(L.north);
    \draw[-{Stealth},acento,line width=1pt] (L.south)--(E.north);
  \end{tikzpicture}
  \end{center}
\end{frame}

% ── CIERRE ──────────────────────────────────────────────────────
{
\setbeamertemplate{background}{
  \begin{tikzpicture}[remember picture,overlay]
    \fill[azuloscuro] (current page.south west) rectangle (current page.north east);
    \foreach \r/\op in {3/0.07,5/0.05,7/0.03}{
      \draw[white,opacity=\op,line width=1pt]
        (current page.center) circle(\r);
    }
    \draw[acento,opacity=0.4,line width=1.5pt]
      (current page.center) circle(2.5);
  \end{tikzpicture}
}
\begin{frame}[plain]
  \begin{center}
    \vspace{1.8cm}
    {\color{acento}\Huge $\vL = I\,\vomega$}\\[0.8cm]
    {\color{white}\large ¡Gracias!}\\[0.3cm]
    {\color{white!60}\small Ricardo · 19 de mayo de 2026}
  \end{center}
\end{frame}
}

\end{document}
